\chapter{Complexity — Straight Line Small}
%%% generated by the blueprint pipeline from TCSlib.Complexity.CircuitComplexity.StraightLineSmall
\section{Overview}

Declarations of the Lean module \texttt{TCSlib.Complexity.CircuitComplexity.StraightLineSmall}.

\section{Declarations}

\begin{definition}[Literal operand]
\lean{BoolCircuit.SLOperand.lit}
\label{BoolCircuit.SLOperand.lit}
\leanok
\uses{BoolCircuit.SLOperand}
For an input index $k\in\mathrm{Fin}(n)$ and a Boolean sign, the corresponding literal
operand is $x_k$ when the sign is true and $\neg x_k$ when it is false.
\end{definition}

\begin{theorem}[Literals precede every program line]
\lean{BoolCircuit.SLOperand.lit_before}
\label{BoolCircuit.SLOperand.lit_before}
\leanok
\uses{BoolCircuit.SLOperand.Before, BoolCircuit.SLOperand.lit}
\difficulty{3}
For any input index $k\text{ Fin}(n)$, Boolean sign $b$, and natural number $i$, the
operand representing the literal that is $x_k$ when $b$ is true and $
eg x_k$ when $b$ is false precedes line $i$.
\end{theorem}

\begin{theorem}[Evaluation of a literal operand]
\lean{BoolCircuit.SLOperand.lit_eval}
\label{BoolCircuit.SLOperand.lit_eval}
\leanok
\uses{BoolCircuit.SLOperand.eval, BoolCircuit.SLOperand.lit}
\difficulty{4}
Let $l=(k,b)$ consist of an input index $k\text{ Fin}(n)$ and a Boolean sign $b$, let
$x:\text{Fin}(n)\to\text{Bool}$ be an assignment, and let $ys$ be a Boolean list. The
literal operand specified by $l$ evaluates to true exactly when $x(k)=b$.
\end{theorem}

\begin{definition}[Three-line circuit shape]
\lean{BoolCircuit.Shape3}
\label{BoolCircuit.Shape3}
\leanok
\uses{BoolCircuit.SLOp}
For a natural number $n$, a three-line circuit shape specifies the operation of each of
its three lines, four signed input literals supplying the first two lines, and the two
operation types conjunction and disjunction available at each line. The third line
combines the values of the first and second lines.
\end{definition}

\begin{definition}[Statements of a three-line shape]
\lean{BoolCircuit.Shape3.stmts}
\label{BoolCircuit.Shape3.stmts}
\leanok
\uses{BoolCircuit.SLOperand.lit, BoolCircuit.SLStmt, BoolCircuit.Shape3}
Let $p$ be a three-line circuit shape on $n$ inputs. Its statement list has three
entries: the first two apply their specified operations to their respective pairs of
signed input literals, and the third applies its specified operation to the values of
lines $0$ and $1$.
\end{definition}

\begin{theorem}[Validity of the three-line program]
\lean{BoolCircuit.Shape3.valid}
\label{BoolCircuit.Shape3.valid}
\leanok
\uses{BoolCircuit.SLOperand.lit_before, BoolCircuit.Shape3, BoolCircuit.Shape3.stmts, BoolCircuit.StraightLine.LinesValid}
\difficulty{4}
For any natural number $n$ and any three-line shape $p$ on $n$ inputs, each of the three
statements in its associated list uses only input literals or values from earlier lines:
the first two statements use input literals, and the third uses the values of the first
and second lines.
\end{theorem}

\begin{definition}[Program of a three-line shape]
\lean{BoolCircuit.Shape3.prog}
\label{BoolCircuit.Shape3.prog}
\leanok
\uses{BoolCircuit.Shape3, BoolCircuit.Shape3.stmts, BoolCircuit.Shape3.valid, BoolCircuit.StraightLineProgram}
Let $p$ be a three-line circuit shape on $n$ inputs. Its straight-line program is the
program whose execution-order statements are those specified by $p$, with each line
using only inputs or earlier lines.
\end{definition}

\begin{theorem}[Length of the three-line program]
\lean{BoolCircuit.Shape3.prog_length}
\label{BoolCircuit.Shape3.prog_length}
\leanok
\uses{BoolCircuit.Shape3, BoolCircuit.Shape3.prog, BoolCircuit.StraightLineProgram.length}
\difficulty{1}
For every three-line shape $p$, the associated straight-line program has length $3$.
\end{theorem}

\begin{theorem}[Evaluation of a three-line program]
\lean{BoolCircuit.Shape3.prog_eval}
\label{BoolCircuit.Shape3.prog_eval}
\leanok
\uses{BoolCircuit.SLOp.apply, BoolCircuit.SLOperand.eval, BoolCircuit.SLOperand.lit_eval, BoolCircuit.SLStmt.eval, BoolCircuit.Shape3, BoolCircuit.Shape3.prog, BoolCircuit.Shape3.prog_length, BoolCircuit.Shape3.stmts, BoolCircuit.StraightLine.runLines, BoolCircuit.StraightLineProgram.eval, BoolCircuit.StraightLineProgram.values}
\difficulty{3}
Let $p$ specify three straight-line Boolean operations, with operations $o_1,o_2,o_3$
and literals $l_1,l_2,l_3,l_4$, and let $x$ be a Boolean assignment to its inputs. The
program evaluates to
$o_3\bigl(o_1(x_{l_1}=b_1,x_{l_2}=b_2),o_2(x_{l_3}=b_3,x_{l_4}=b_4)\bigr)$, where each
literal $l_i$ consists of an input coordinate and a Boolean sign $b_i$, and
$x_{l_i}=b_i$ denotes equality of that input value with the sign.
\end{theorem}

\begin{definition}[Single-line straight-line program]
\lean{BoolCircuit.oneLine}
\label{BoolCircuit.oneLine}
\leanok
\uses{BoolCircuit.SLOp, BoolCircuit.SLOperand.lit, BoolCircuit.SLOperand.lit_before, BoolCircuit.StraightLineProgram}
For a natural number $n$, an operation that is either conjunction or disjunction, and
two signed input literals, the corresponding straight-line program consists of one
statement applying that operation to those literals.
\end{definition}

\begin{theorem}[Length of a single-gate program]
\lean{BoolCircuit.oneLine_length}
\label{BoolCircuit.oneLine_length}
\leanok
\uses{BoolCircuit.SLOp, BoolCircuit.StraightLineProgram.length, BoolCircuit.oneLine}
\difficulty{1}
Let $n$ be a natural number, let $o$ be conjunction or disjunction, and let $l_1,l_2$ be
signed input literals among the $n$ inputs. The straight-line program consisting of the
single gate $o$ applied to $l_1$ and $l_2$ has length $1$.
\end{theorem}

\begin{theorem}[Evaluation of a single gate]
\lean{BoolCircuit.oneLine_eval}
\label{BoolCircuit.oneLine_eval}
\leanok
\uses{BoolCircuit.SLOp, BoolCircuit.SLOp.apply, BoolCircuit.SLOperand.lit_eval, BoolCircuit.SLStmt.eval, BoolCircuit.StraightLine.runLines, BoolCircuit.StraightLineProgram.eval, BoolCircuit.StraightLineProgram.values, BoolCircuit.oneLine, BoolCircuit.oneLine_length}
\difficulty{3}
Let $n$ be a natural number, let $o$ be either conjunction or disjunction, let
$l_1,l_2\in\{0,\ldots,n-1\}\times\{\mathrm{false},\mathrm{true}\}$ be signed input
literals, and let $x$ be a Boolean assignment to the $n$ inputs. The program consisting
of the single gate $o$ applied to $l_1,l_2$ evaluates to $o$ applied to the truth values
of the literals: the value at index $l_i$ agrees with the sign of $l_i$ for each
$i=1,2$.
\end{theorem}

\begin{definition}[Boolean Gate Lookup Table]
\lean{BoolCircuit.table1}
\label{BoolCircuit.table1}
\leanok
\uses{BoolCircuit.SLOp}
Define a table on pairs of Boolean inputs whose output consists of a conjunction or
disjunction and two references to the sole position $0$, each tagged by a Boolean. On
inputs $(\text{false},\text{false})$ it returns conjunction with tags
$(\text{true},\text{false})$; on $(\text{false},\text{true})$ it returns conjunction with
tags $(\text{true},\text{true})$; on $(\text{true},\text{false})$ it returns conjunction
with tags $(\text{false},\text{false})$; and on $(\text{true},\text{true})$ it returns
disjunction with tags $(\text{true},\text{false})$.
\end{definition}

\begin{definition}[Two-Step Boolean Lookup Table]
\lean{BoolCircuit.table2}
\label{BoolCircuit.table2}
\leanok
\uses{BoolCircuit.Shape3}
For four Boolean inputs $(a,b,c,d)$, define a lookup table returning three gate
operations and four signed references to vertices $0$ and $1$. In the order
$(a,b,c,d)=0000,0001,0010,0011,0100,0101,0110,0111,1000,1001,1010,1011,1100,1101,1110,1111$,
its outputs are respectively: $(\land;(0,+),(0,+);\land;(0,+),(0,-);\land)$;
$(\land;(0,+),(0,+);\land;(0,+),(1,+);\land)$;
$(\land;(0,-),(0,-);\land;(0,-),(1,+);\land)$;
$(\land;(1,+),(1,+);\land;(1,+),(1,+);\land)$;
$(\land;(0,+),(0,+);\land;(0,+),(1,-);\land)$;
$(\land;(0,+),(0,+);\land;(0,+),(0,+);\land)$;
$(\land;(0,+),(1,-);\land;(0,-),(1,+);\lor)$;
$(\land;(0,+),(0,+);\land;(0,-),(1,+);\lor)$;
$(\land;(0,-),(0,-);\land;(0,-),(1,-);\land)$;
$(\land;(0,+),(1,+);\land;(0,-),(1,-);\lor)$;
$(\land;(0,-),(0,-);\land;(0,-),(0,-);\land)$;
$(\land;(0,+),(1,+);\land;(0,-),(0,-);\lor)$;
$(\land;(1,-),(1,-);\land;(1,-),(1,-);\land)$;
$(\land;(0,+),(0,+);\land;(0,-),(1,-);\lor)$;
$(\land;(0,+),(1,-);\land;(0,-),(0,-);\lor)$;
$(\land;(0,+),(0,+);\land;(0,-),(0,-);\lor)$. Here each signed reference $(i,+)$ or
$(i,-)$ refers to vertex $i$ or its Boolean negation, and the three operations are those
of the first, second, and third lines.
\end{definition}

\begin{theorem}[One-line representation of a Boolean function]
\lean{BoolCircuit.exists_oneLine}
\label{BoolCircuit.exists_oneLine}
\leanok
\uses{BoolCircuit.SLOp.apply, BoolCircuit.StraightLineProgram, BoolCircuit.StraightLineProgram.eval, BoolCircuit.StraightLineProgram.length, BoolCircuit.oneLine, BoolCircuit.oneLine_eval, BoolCircuit.table1}
\difficulty{5}
Every Boolean function on one input bit is computed by a valid straight-line program of
length exactly one.
\end{theorem}

\begin{theorem}[Three-line representation of a Boolean function]
\lean{BoolCircuit.exists_threeLine}
\label{BoolCircuit.exists_threeLine}
\leanok
\uses{BoolCircuit.SLOp.apply, BoolCircuit.Shape3.prog, BoolCircuit.Shape3.prog_eval, BoolCircuit.StraightLineProgram, BoolCircuit.StraightLineProgram.eval, BoolCircuit.StraightLineProgram.length, BoolCircuit.table2}
\difficulty{4}
Every Boolean function on two input bits is computed by a valid straight-line program of
length exactly three.
\end{theorem}

\begin{theorem}[Straight-line simulation for at most two inputs]
\lean{BoolCircuit.exists_straightLine_length_le_size_of_le_two}
\label{BoolCircuit.exists_straightLine_length_le_size_of_le_two}
\leanok
\uses{BoolCircuit.DAGCircuit, BoolCircuit.DAGCircuit.eval, BoolCircuit.DAGCircuit.size, BoolCircuit.DAGCircuit.values_getD_input, BoolCircuit.StraightLineProgram, BoolCircuit.StraightLineProgram.eval, BoolCircuit.StraightLineProgram.length, BoolCircuit.exists_oneLine, BoolCircuit.exists_threeLine, BoolCircuit.oneLine, BoolCircuit.oneLine_eval, BoolCircuit.oneLine_length}
\difficulty{5}
Let $n$ satisfy $1\le n\le 2$, and let $C$ be a directed acyclic Boolean circuit on $n$
inputs. There is a valid straight-line program on the same inputs, of length at most the
size of $C$, that computes the same Boolean function as $C$.
\end{theorem}
